\documentclass[10pt,a4paper]{amsart}
\usepackage[margin=19mm]{geometry}
\usepackage{amsmath,amssymb,mathtools}
\usepackage[scr=rsfs,cal=euler]{mathalfa}
\usepackage{xfrac}
\usepackage[unicode,hidelinks,bookmarks=false]{hyperref}
\newcommand{\ie}{i.e.\ }
\newcommand{\cf}{cf.\ }
\newcommand{\resp}{respectively\ }
\newcommand{\Subsection}{\subsection}
\providecommand{\nobreakdash}{\nobreak\hbox{-}}
\setlength{\emergencystretch}{2em}
\multlinegap0pt
\usepackage{amssymb}
\usepackage{bm}
\usepackage{xcolor}
\usepackage{tikz}
\usepackage[shortlabels]{enumitem}
\usepackage[normalem]{ulem}
\usepackage{mathtools}
\newtheorem{proposition}{Proposition}
\newtheorem{theorem}{Theorem}
\usepackage{cleveref}
\crefname{proposition}{Proposition}{Propositions}
\crefname{theorem}{Theorem}{Theorems}
\newcommand{\CC}{\mathbb{C}}
\title{On Waring Rank Jumps \\ via critical rank-one approximations}
\author{Alessandro Oneto, Pierpaola Santarsiero, Ettore Teixeira Turatti}
\date{Source: arXiv:2605.04695v1 (CC BY 4.0)}
\begin{document}
\maketitle

\noindent\emph{Source and licence.} On Waring Rank Jumps \\ via critical rank-one approximations. Version arXiv:2605.04695v1 (CC BY 4.0), distributed by arXiv under the Creative Commons Attribution 4.0 International licence, http://creativecommons.org/licenses/by/4.0/, as recorded in the arXiv licence metadata for this exact version and shown on the abstract page https://arxiv.org/abs/2605.04695v1. Full licence notice: \texttt{licenses/arxiv-2605.04695-CC-BY-4.0.txt}.\par\smallskip

\noindent\textit{Editorial scope.} Counted unit: the article's theorem rnc (source0.tex lines 811--813). The article's complete original proof of it is delivered with it (source0.tex lines 814--855, one \texttt{proof} environment of 42 lines). No mathematics is written, repaired or re-derived: the statement and the proof are the article's own lines, in the article's own notation, and no retained line is substituted or re-flowed.  Retained uncounted as the source-local context the proof needs: lines 681--683; lines 737--739.  Nothing was omitted from any retained span, so the package carries no omission record.  No retained line contains a citation, so the wrapper carries no bibliography and no bibliography entry.  No retained line contains a figure environment, an image-inclusion command or drawing code, so no image is delivered and the package declares no asset dependency.  Editorial formatting only, in addition: an AMS \texttt{amsart} preamble that restates the article's own declarations and macros used by the retained lines, namely the environments \texttt{proposition}, \texttt{theorem} on separate counters, together with the standard packages those lines need.  The retained environments print in this document as Proposition 1, Theorem 1 in order of appearance, whereas the article prints the same items with the numbering of its own full document.  The complete native source of the article as distributed in the arXiv e-print for this exact version is bundled unchanged as \texttt{sources/199885/source0.tex}.
     \begin{align}\label{eq: equations gj}
     g_j=\sum_{i=1}^{r-1}\alpha_{i,0}^{d-1}\alpha_{i,j} \text{ for }j=1,\dots, n
     \end{align}

\begin{proposition}\label{prop:rankrbinary}
    The variety defined by the equations $g_1,\dots,g_n$ of \cref{eq: equations gj} contains a polynomial of Waring rank $r\leq \frac{d+1}{2}$.
\end{proposition}

\begin{theorem}\label{proposition: dim WE rnc }
    Let $2\leq r\leq \frac{d+1}{2}$, then $\dim(\mathrm{WE}_{1,r}^d)=2(r-1)$. In particular, $\mathrm{WE}_{1,r}^d$ is codimension one subvariety of $\sigma_r(V_1^d)$.
\end{theorem}

\begin{proof}
    {The strategy is: we express $\mathrm{WE}^d_{1,r}$ as the image of a map 
    $
        \psi:\mathrm{O}(2)\times Y\to \sigma_r(V_1^d),\ \psi(g,y)=g\cdot y,
    $
    then, we identify a specific element for which the fiber $\psi^{-1}(z)$ is zero-dimensional and conclude by semi-continuity.}
    
    Since $\sigma_{r-1}(V_1^d)\subsetneq\mathcal J([x_0^d], \sigma_{r-1}(V_1^d))$, then $\dim \mathcal J([x_0^d], \sigma_{r-1}(V_1^d)) = \dim \sigma_{r-1}(V_1^d)+1 = 2(r-1).$ Let $Y=H_1\cap \mathcal J([x_0^d], \sigma_{r-1}(V_1^d))$ We have $\dim(Y)=2(r-1)-1$.
    As mentioned, $\mathrm{WE}_{1,r}^d=\mathrm{im}(\psi)$. Thus, $$\dim \mathrm{WE}_{n,r}^d=\dim Y+\dim \mathrm{O}(2)-\dim \psi^{-1}(z),$$ where $z$ is a generic element in the image of $\psi$ and $\psi^{-1}(z)=\{(g,y)\in \mathrm{O}(2)\times Y\mid g\cdot y=z\}$. Let $\pi_2:\mathrm{O}(2)\times Y\to Y$ be the projection onto the second factor, so $$\dim \psi^{-1}(z)=\dim \pi_2(\psi^{-1}(z))+\dim \pi_2^{-1}(y)$$ for a generic $y\in \pi_2(\psi^{-1}(z))$. Notice $\pi_2(\psi^{-1}(z))=\{y\mid g\cdot y=z \text{ for some $g\in\mathrm{O}(2)$}\}$ and $\pi_2^{-1}(y)=\{(g,y)\mid g\cdot y=z\}$. We will show that both are zero-dimensional for a specific choice of $z$, and by semi-continuity, also the generic fibre of $\psi$ is zero-dimensional. 
    
    We fix $z=x_0^d+\sum_{i=1}^{r-1}(x_0+\xi^ix_1)^d\in Y\subset \mathrm{im}(\psi)$, for $\xi$ a fundamental $(r-1)$-root of unity, where the belonging follows from the proof of \Cref{prop:rankrbinary}.

    Since $y\in Y$, $y=x_0^d+\sum_{i=1}^{r-1}L_i^d$, and notice that $y\in \pi_2(\psi^{-1}(z))$ if and only if $y\in Y\cap \mathrm{O}(2)\cdot z$. By assumption, both $y$ and $z$ have rank $r\leq \frac{d+1}{2}$, both are identifiable, therefore $g\cdot\{ x_0,(x_0+\xi^ix_1)\}_{i=1}^{r-1}=\{\mu_0 x_0,\mu_iL_i\}_{i=1}^{r-1}$, where $\mu_0^d=\mu_i^d=1$. It follows that either $g\cdot x_0=\mu_0 x_0$ or $g\cdot (x_0+\xi^ix_1)=\mu_0x_0$ for some $i$. 

    If $g\cdot x_0=\mu_0 x_0$, then $g=\begin{bmatrix}
        \mu_0 &0\\
        0&\eta
    \end{bmatrix}$, with $\mu_0^2=\eta^2=1$, and we obtain finitely many such $y$'s. On the other hand, if $g\cdot \frac{1}{\|(1,\xi^i)\|}(x_0+\xi^ix_1)=\mu_0 x_0$, then 
    $$
    g=\frac{1}{\|(1,\xi^i)\|}\begin{bmatrix}
       \mu_0-\xi^{2i}&\xi^i\\
       -\xi^i&\mu_0-\xi^{2i}
    \end{bmatrix} \quad \text{ or } \quad g=\frac{1}{\|(1,\xi^i)\|}\begin{bmatrix}
       \mu_0-\xi^{2i}&\xi^i\\
       \xi^i&-\mu_0+\xi^{2i}
    \end{bmatrix}.
    $$
   It follows that for each $i=1,\dots, r-1$, there are at most finitely many such $y$'s in $Y\cap \mathrm{O}(2)\cdot z$. Finally, we notice that if $\|(1,\xi^i)\|=0$, i.e., it is isotropic, then there is no $g\in \mathrm{O}(2)$ such that $g\cdot x_0=\lambda (x_0+\xi^ix_1)$ for any $\lambda\in \CC$, and the finiteness argument still holds. Thus $\dim \pi_2(\psi^{-1}(z))=0$.

    We now focus on $\pi_2^{-1}(y)$: let $y=x_0^d+\sum_{i=1}^{r-1}L_i^d\in \pi_2(\psi^{-1}(z))$, since $g\cdot y=z$, and $y$ is identifiable, we have $g\cdot \{x_0,L_i\}_{i=1}^{r-1}=\{\mu_0x_0,\mu_i(x_1+\xi^ix_1)\}_{i=1}^{r-1}$, with $\mu_0^d=\mu_i^d=1$. Similarly, either $g\cdot x_0=\mu_0x_0$ or $g\cdot x_0=\frac{\mu_i}{\|(1,\xi^i)\|}(x_0+\xi^ix_1)$. For the first, again it follows that $g$ is diagonal with entries $\pm 1$. For the second, it follows that
  $$
 g=\frac{1}{\|(1,\xi^i)\|}\begin{bmatrix}
        1&\xi^i\\
        -\xi^i&1
    \end{bmatrix} \quad \text{ or }  \quad g=\frac{1}{\|(1,\xi^i)\|}\begin{bmatrix}
        1&\xi^i\\
        \xi^i&-1
    \end{bmatrix} ,$$  
    and there are finally many possibilities. Again, if $(1,\xi^i)$ is isotropic, there are no such matrices, and finiteness still holds.

   Hence $\psi^{-1}(z)$ is zero-dimensional, so also the generic fibre is zero-dimensional, and we conclude that $\dim(\mathrm{WE}_{1,r}^d)=\dim Y+\dim \mathrm{O}(2)=2(r-1)$.
    \end{proof}

\end{document}
